\begin{ledgroupsized}[r]{120mm}
\footnotesize
\pstart
\noindent\textbf{Überlieferung:}
\pend
\end{ledgroupsized}
\begin{ledgroupsized}[r]{114mm}
\footnotesize
\pstart
\parindent -6mm
\makebox[6mm][l]{\textit{L}}%
Konzept:
LH~XXXV~10, 14~Bl.~8. 
Ein Blatt~4\textsuperscript{o};
\href{\Scan{LH_35_10_14_008wz}}%
{Wasserzeichenfragment} am Blattrand;
oberer und linker Rand beschnitten.
Eine Seite auf Bl.~8~r\textsuperscript{o}; 
Bl.~\href{\Scan{LH_35_10_14_008v}}%
{\lbrack8~v\textsuperscript{o}\rbrack}
leer bis auf zwei gestrichene und einen nicht gestrichenen Entwurf zu \bracks{\textit{Fig.~1}}, die nicht wiedergegeben werden.
Am oberen Rand von Bl.~8~r\textsuperscript{o}, in Rot, der Vermerk von unklarer Hand \textit{VIII}
(vgl. die Träger der benachbarten Signaturen LH~XXXV~10, 14 Bl.~2 bis 13, u.\,a. 
N.\;\ref{60117},
\ref{60118},
\ref{60112},
\ref{60116},
\ref{60113},
\ref{60114},
\ref{60110},
\ref{41824}
und \ref{60111}
in diesem Band).
\pend
\end{ledgroupsized}
\vspace{5mm}
\begin{ledgroup}
\footnotesize
\pstart
\noindent
\textbf{Datierungsgründe:}
Das vorliegende Konzept N.\;\ref{60115} hat Leibniz auf einem Blatt verfasst, dessen Wasserzeichenfragment einerseits bislang nicht identifiziert werden konnte, das andererseits ähnliche Eigenschaften wie andere Papiere italienischer oder österreichischer Herkunft aus der Zeit der Italienreise (1688--1690) aufweist.
Maßgebend für die angegebene Datierung ist allerdings nicht die (unsichere) Bestimmung des Trägers, sondern folgender inhaltlicher Umstand.
\pend
\pstart
Leibniz beruft sich in N.\;\ref{60115} zur Beleuchtung einer statischen Frage auf einen Fall der Stoßlehre, der eine angemessene Analogie für die Statik bietet, mit den Worten
"Confugiendum est ad ea quae de directione Eccentrica a me primum sunt definita" (\refpassage{35_10_14_008-1a}{35_10_14_008-1b}).
Gemeint sind die Ergebnisse der intensiven Beschäftigung mit dem exzentrischen Stoß am Ende der italienischen Zeit, die zu innovativen Ergebnissen führte (in den Worten von N.\;\ref{60115}: "\textit{a me primum} sunt definita").
Sämtliche erhaltene Konzepte sind nach der Abreise aus Florenz\indexO{Florenz}, während der Rückreise nach Deutschland oder in den ersten Monaten in Hannover\indexO{Hannover} entstanden:
die frühesten (N.\;\ref{impulsuseccentricus}) ab Dezember 1689,
die letzten (N.\;\ref{cecs_intro}) im September 1691.
Zwar hatte Leibniz beabsichtigt, im letzten Abschnitt der \cite{01354}\textit{Dynamica de potentia}, "De Concursu Corporum", auch den exzentrischen Stoß zu behandeln, aber er konnte das Werk in Florenz nicht vollenden: Der Abschnitt zur Stoßlehre bricht vor dieser geplanten Besprechung ab
(\textit{CFP} I, S.\;444, Fassung \bracks{A}).
\pend
\pstart
Die frühesten erhaltenen Konzepte zu diesem Thema (N.\;\ref{impulsuseccentricus}, bes. N.\;\ref{60142+60297}: "Impulsus eccentrice directi effectum explicare") folgen noch der Form "Propositio ... problema", in welcher Leibniz sie in die Beweisstruktur der \textit{Dynamica} zu integrieren trachtete (zur Entwicklung der Leibniz'schen Nachforschungen über den exzentrischen Stoß nach 1689 siehe die Datierung von N.\;\ref{impulsuseccentricus}).
Eine weitere Anspielung auf Befunde zum "nisum gyrandi" ("Manifestum est ex alibi explicatis...") im selben Absatz von N.\;\ref{60115} gilt tatsächlich dem erwähnten Text N.\;\ref{60142+60297}, denn Leibniz beruft sich ausdrücklich auf die dort aufgestellte These, dass dem direkten exzentrischen Stoß eine Drehung um den Mittelpunkt des indirekt angestoßenen Körperteils folge (siehe dort \refpassage{60142+60297_centrumgyrationis-1}{60142+60297_centrumgyrationis-2} und \refpassage{LH_37_05_043_043v_centrgyr-1}{LH_37_05_043_043v_urzelle-2}).
Diese Theorie verwirft Leibniz allerdings wenige Monate später endgültig, zuerst im Konzept "De impulso eccentrico" von April -- Oktober 1690 (N.\;\ref{60300}).
N.\;\ref{60142+60297} kann frühestens im Dezember 1689 entstanden sein, während  N.\;\ref{60300} spätestens im Oktober 1690 verfasst wurde; die Abfassung unseres Konzepts N.\;\ref{60115} muss in der Zeitspanne zwischen diesen beiden Endpunkten fallen.
\pend 
\end{ledgroup}

\clearpage
\selectlanguage{latin}		%
\vspace{8mm}
\pstart
\noindent
\href{\Scan{LH_35_10_14_008r}}%
{\lbrack8~r\textsuperscript{o}\rbrack} 
\pend
\nobreak
\begin{edfigure}[0em][1em]
\hfill
\includegraphics[scale=1]{%
gesamttex/edit_VIII-4/images/LH_35_10_14_008_d1_008r.pdf}
\hfill%
\includegraphics[scale=1]{%
gesamttex/edit_VIII-4/images/LH_35_10_14_008_d2_008r.pdf}
\hfill%
\edcaption{%
\edtext{%
\hspace*{18mm}%
\bracks{\textit{Fig.~1}}}{%
\lemma{\bracks{\textit{Fig.~1}}}%
\Cfootnote[nonum]{%
Zwei gestrichene und ein nicht gestrichener Entwurf zum Diagramm auf Bl.~8~v\textsuperscript{o} werden nicht wiedergegeben.}}
\hspace*{45mm}%
\bracks{\textit{Fig.~2}}}
\end{edfigure}
\pstart
\noindent
Sit globus\indexS{globus} 
\edtext{\textit{ABC}}{%
\lemma{}%
\Bfootnote{%
\textit{ABC} \textit{erg.~L}}}
ex materia similari constans, seu centrum saltem gravitatis\indexS{centrum gravitatis} \textit{C} habens in ipso centro globi.\indexS{centrum globi} Is sustineatur\indexS{globus a planis inclinatis sustentatus} si 
\edtext{placet in punctis}{%
\lemma{placet}%
\Bfootnote{%
\textit{(1)}~a duobus %
\textit{(2)}~in punctis~\textit{L}}}
\textit{A}, \textit{B} a planis \textit{LM}, \textit{NP}, quaeritur qua 
\edtext{vi plana\indexS{planum inclinatum}}{%
\lemma{vi}%
\Bfootnote{%
\textit{(1)}~pondere %
\textit{(2)}~plana~\textit{L}}}
premantur. Vel pro planis, sint sustentacula\indexS{sustentaculum} seu pali\indexS{palus} \textit{AE}, \textit{BF}, ad dicta plana\indexS{planum inclinatum} perpendiculares adeoque producti concurrentes\indexS{directiones sustentaculorum in centro concurrentes} in centro\indexS{centrum gravitatis} \textit{C}, 
\edtext{quaeritur quantam}{%
\lemma{quaeritur}%
\Bfootnote{%
\textit{(1)}~quae vis sit sustenta %
\textit{(2)}~quantam~\textit{L}}}
sustentacula\indexS{sustentaculum} vim\indexS{vis quam patitur sustentaculum} patiantur. Seu si ponamus palos\indexS{palus} \textit{AE}, \textit{BF} esse normales 
\edtext{ad radios\indexS{radius mobilis circa centrum} \textit{GE}, \textit{HF}}{%
\lemma{ad}%
\Bfootnote{%
\textit{(1)}~radiu %
\textit{(2)}~radios \textit{GE}, %
\textit{(a)}~\textit{FH}, circa cent %
\textit{(b)}~\textit{HF}~\textit{L}}}
circa centra \textit{G} et \textit{H} mobiles, esseque circa idem centrum mobiles\indexS{radius mobilis circa centrum} alios radios horizonti\indexS{horizon} parallelos \textit{GI}, \textit{HK} quorum extremis \textit{I} et \textit{K} appensa sint pondera, quaeritur, quam proportionem debeant habere pondera \textit{I} et \textit{K}, inter se, et cum pondere \textit{C}, ut omnia sint in aequilibrio.\indexS{aequilibrium} 
Quoniam igitur perinde est ac 
\edtext{si omnis}{%
\lemma{si}%
\Bfootnote{%
\textit{(1)}~punctum %
\textit{(2)}~omnis~\textit{L}}}
gravitas collecta esset in punctum \textit{C}, idque punctum nisibus 
\edtext{in \textit{AC},}{%
\lemma{in}%
\Bfootnote{%
\textit{(1)}~\textit{EC}, %
\textit{(2)}~\textit{AC},~\textit{L}}}
\textit{BC}, sustineretur. Patet nisus\indexS{nisus} esse tales, ut inter se compositi\indexS{compositio nisuum} aequent directionem 
\edtext{gravis \textit{C}}{%
\lemma{gravis}%
\Bfootnote{%
\textit{(1)}~horizonta ad %
\textit{(2)}~\textit{C}~\textit{L}}}
ad horizontem. Itaque sumta \textit{CQ}, 
\edtext{pro nisu gravis,\indexS{nisus gravis} in linea directionis\indexS{linea directionis}}{%
\lemma{pro}%
\Bfootnote{%
\textit{(1)}~directione gravis %
\textit{(2)}~nisu gravis in linea directionis~\textit{L}}}
seu perpendiculari ad horizontem; et ex \textit{Q} educendo
\edtext{ipsis directionibus sustentaculorum\indexS{sustentaculum} seu}{%
\lemma{ipsis}%
\Bfootnote{%
\textit{(1)}~directionis
\textit{(2)}~directionibus 
\textit{(a)}~in
\textit{(b)}~sustentaculorum
\textit{(aa)}~parallelas %
\textit{(bb)}~seu~\textit{L}}}
ipsis \textit{AC}, \textit{BC}, parallelas \textit{QR}, \textit{QS}; patet fore pondera \textit{I}, \textit{C}, \textit{K}, ut rectas \textit{CS}, \textit{CQ}, \textit{CR}.
\pend
\pstart
Sed quid 
\edtext{si directiones}{%
\lemma{si}%
\Bfootnote{%
\textit{(1)}~puncta %
\textit{(2)}~directiones~\textit{L}}}
sustinentium\indexS{directio sustinentis} non tendant\indexS{directiones sustinentium non in centrum tendentes} in centrum gravitatis;\indexS{centrum gravitatis} vel saltem non omnes, adeoque nec concurrant inter se\bracks{?}\indexS{directiones inter se non concurrentes} 
\edlabel{35_10_14_008-1a}%
Confugiendum est ad 
\edtext{ea quae de directione Eccentrica\indexS{directio eccentrica} a me primum sunt definita.\edlabel{35_10_14_008-1b}}{%
\lemma{ea \bracks{...} definita}%
\Cfootnote{Gemeint sind die Befunde zum exzentrischen Stoß in nach der italienischen Zeit entstandenen Konzepten (siehe die "Datierungsgründe"), insbesondere N.\;\ref{60142+60297}: siehe die Erläuterung zu \refnote{FnC:LH_35_10_14_008-gyrandi}.}}
Ponamus ergo mobilis centrum\indexS{centrum gravitatis} esse quidem \textit{C}. Sed 
\edtext{unius}{%
\lemma{}%
\Bfootnote{%
unius \textit{erg.~L}}}
sustentaculi directionem\indexS{directio sustentaculi} esse \textit{XT} tendentem in \textit{T} centrum portionis ipsius mobilis \textit{TCV}, et reliquae portionis centrum esse \textit{V}.
\edtext{Manifestum est ex alibi explicatis perinde esse ac si mobile acciperet nisum gyrandi\indexS{nisus gyrandi}}{%
\applabel{FnC:LH_35_10_14_008-gyrandi}%
\lemma{Manifestum \bracks{...} gyrandi}%
\Cfootnote{Gemeint ist das ursprünglich wohl für den Schlussteil der \cite{01354}\textit{Dynamica de potentia} vorgesehene Konzept N.\;\ref{60142+60297}.}}  
circa \textit{V}, sic
\edtext{ut celeritas\indexS{celeritas} vel (si in mortuis nisibus\indexS{nisus mortuus} sic appellare malis) promtitudo\indexS{promtitudo nisus}}{%
\lemma{ut}%
\Bfootnote{%
\textit{(1)}~celeritas %
\textit{(2)}~prom
\textit{(3)}~celeritas 
\textit{(a)}~(vel in mortuis cona
\textit{(a)}~vel (si \bracks{...} promtitudo~\textit{L}}}
ipsius \textit{T} 
\edtext{respondeat promtitudini nisus\indexS{promtitudo nisus} qui}{%
\lemma{respondeat}%
\Bfootnote{%
\textit{(1)}~nisum %
\textit{(2)}~nisui \textit{XT} %
\textit{(3)}~promtitudini nisus \bars{\textit{XT} \textit{erg. u. gestr.}} %
\textit{(a)}~se %
\textit{(b)}~qui~\textit{L}}}
est in directione \textit{XT}, unde promtitudo\indexS{promtitudo nisus} ipsius \textit{C} erit ad promtitudinem\indexS{promtitudo nisus} ipsius \textit{T}, seu nisus sustentaculi,\indexS{nisus sustentaculi} ut \textit{CV} 
\edtext{ad \textit{TV}. Alterius}{%
\lemma{ad \textit{TV}.}%
\Bfootnote{%
\textit{(1)}~Ipsius %
\textit{(2)}~Alterius~\textit{L}}}
autem sustentaculi nisus\indexS{nisus sustentaculi} secundum \textit{YC} ponatur tendere\indexS{nisus in centrum gravitatis tendens} in ipsum centrum gravitatis mobilis.\indexS{centrum gravitatis mobilis} Sumta igitur \textit{CZ} in linea directionis ipsius \textit{C}, ducantur ipsae \textit{Z$\alpha$}, \textit{Z$\beta$}, parallelae ipsis \textit{XT}, \textit{YC} et compleatur parallelogrammum \textit{C$\alpha$Z$\beta$}, tunc 
\edtext{ipsa \textit{CZ}}{%
\lemma{ipsa}%
\Bfootnote{%
\textit{(1)}~\textit{C$\alpha$} %
\textit{(2)}~\textit{CZ}~\textit{L}}}
exprimente totum ponderis \textit{C} nisum\bracks{,}\indexS{nisus ponderis} 
ipsa \textit{C$\alpha$} exprimet nisum sustentaculi\indexS{nisus sustentaculi}
\textit{Y}. Sed nisus sustentaculi\indexS{nisus sustentaculi}
\textit{X}, erit ad ipsam \textit{C$\beta$}, in ratione composita \textit{TV} ad \textit{CV}, et gravis \textit{T} (portionis) ad grave \textit{C} (totum).
\pend
